% !TEX program = xelatex
%==============================================================================
%  2026 五一数学建模竞赛（校赛模拟）B 题 —— 主文档
%  模板：gmcmthesis.cls（vendored，未改动；源自 zhanwen/MathModel 2025 模版）
%  引擎：xelatex + biber    构建：make   （或 latexmk -xelatex main.tex）
%  规范：见同目录 TYPESPEC.md
%==============================================================================
\documentclass{gmcmthesis}

\usepackage{amsmath}
\usepackage{pdfpages}
\usepackage{float}
\graphicspath{{figures/}}

% ---- 参考文献：biblatex + GB/T 7714-2015（顺序编码制）----
\usepackage[backend=biber, style=gb7714-2015, gbnamefmt=lowercase]{biblatex}
\addbibresource{ref.bib}

% ---- 图表按节编号（如 X-1）----
\numberwithin{figure}{section}
\numberwithin{table}{section}
\renewcommand{\thefigure}{\arabic{section}-\arabic{figure}}
\renewcommand{\thetable}{\arabic{section}-\arabic{table}}

% ---- 页眉页脚（覆盖 cls 的 \pagestyle{plain}）----
\usepackage{fancyhdr}
\geometry{headheight=14pt}
\renewcommand{\sectionmark}[1]{\markboth{#1}{}}
\pagestyle{fancy}
\fancyhf{}
\fancyhead[L]{\zihao{-5}\leftmark}                            % 左：当前章节名
\fancyhead[R]{\zihao{-5}五一数学建模竞赛（校赛模拟）· B 题}    % 右：赛题
\fancyfoot[C]{\zihao{-5}-\ \thepage\ -}                        % 中：页码
\renewcommand{\headrulewidth}{0.4pt}
\renewcommand{\footrulewidth}{0pt}

% ---- 论文元信息（团队填写）----
\title{多工序协同作业调度与设备购置优化模型}
\baominghao{16组}
\schoolname{哈尔滨理工大学}
\membera{骆宏博}
\memberb{梁震宇}
\memberc{赵晨曦}

%==============================================================================
%  封面：本次无官方封面，采用「题目版式」封面（仿 cls 原生 \makenametitle：
%  顶部竞赛名文字 + 「题　目」行 + 队员信息表，单页、无图片依赖）。
%==============================================================================
\makeatletter
\renewcommand{\maketitle}{%
  \thispagestyle{empty}%
  \vspace*{2.5cm}
  \begin{center}
    {\xinwei\zihao{2} 数学建模竞赛（校赛模拟）\par}
    \vskip0.4cm
    {\zihao{3}\heiti B\quad 题\par}
    \vskip1.8cm
    \begin{tabularx}{.95\textwidth}{@{}lX@{}}
      {\zihao{4} 题\qquad 目} & \hfill\zihao{-3}\heiti \@title \hfill\null \\
      \cline{2-2}
    \end{tabularx}
  \end{center}
  \vfill
  \renewcommand{\arraystretch}{1.8}%
  \noindent\begin{tabularx}{\textwidth}{lX}
    {\zihao{4}\bfseries 参赛队号} & {\zihao{4} \gmcm@tokens@baominghao} \\
    \cline{2-2}\\[-12pt]
    {\zihao{4}\bfseries 学校名称} & {\zihao{4} \gmcm@tokens@schoolname} \\
    \cline{2-2}\\[-12pt]
    {\zihao{4}\bfseries 队员姓名} & {\zihao{4} 1.\;\gmcm@tokens@membera \quad 2.\;\gmcm@tokens@memberb \quad 3.\;\gmcm@tokens@memberc} \\
    \cline{2-2}
  \end{tabularx}
  \vspace{2.5cm}
  \newpage
}
\makeatother

\begin{document}

\maketitle

%==============================================================================
%  摘要（独立成页；标题「摘　要」由 cls 以隶书渲染）
%  写作要点：方法 → 各问一句话说清思路与关键结果（含数字）→ 创新点。
%  关键词 3–5 个，分号分隔，末尾无标点。
%==============================================================================
\begin{abstract}
针对五个车间集中整修任务中的多工序协同作业问题，本文将其抽象为带车间内工序链约束、同类设备可替代、双设备协同作业以及跨车间序列相关运输时间的多资源调度问题。首先对 C 车间循环工序进行展开，将原始工序转化为 27 个工序实例和 41 个设备子任务；然后建立以最小化全部任务完工时间为目标的混合整数规划模型，并通过设备分配变量、工序开始时间变量和同设备排序变量刻画工序顺序、设备互斥与跨车间转运约束。针对追加设备购置预算的问题，进一步引入采购变量和预算约束，形成采购--调度联合优化模型。为增强结果可信度，本文还使用 CP-SAT 构造问题二的瓶颈资源松弛下界和问题四的显式采购验证模型，对混合整数规划结果进行交叉验证。

求解结果表明：仅由班组 1 完成 A 车间任务时，最短完工时间为 $41600\,\mathrm{s}$，即 11:33:20；仅使用班组 1 完成五个车间任务时，最短完工时间为 $163764\,\mathrm{s}$，即 45:29:24；两个班组共同作业时，最短完工时间缩短为 $123844\,\mathrm{s}$，即 34:24:04；在 50 万元购置预算下，最短完工时间仍为 $123844\,\mathrm{s}$，最优采购方案为不购置任何设备。进一步地，问题一由 A 车间工序链下界证明最优，问题二由高速抛光机和自动传感多功能机双瓶颈资源松弛下界证明在本文双设备同步协同模型下最优，问题三和问题四由 C 车间关键工序链下界及显式采购模型证明最优。该方案不仅给出满足题目要求的设备调度表和采购方案，也为多车间协同调度问题提供了可复现、可解释的精确建模框架。

\keywords{多工序协同作业；柔性作业车间调度；混合整数规划；CP-SAT；设备购置优化}
\end{abstract}

\tableofcontents
\clearpage

% 01_restatement.tex —— 问题重述
\section{问题重述}

某工业制造系统需要对 A、B、C、D、E 五个车间开展集中整修任务。每个车间包含若干具有固定先后顺序的工序，工序编号由小到大表示执行顺序。不同工序需要占用不同类型的设备，部分工序需要两类设备共同参与；若一道工序需要两类设备，则两类设备均需完成该工序对应工程量，只有当两类设备均完成后，该工序才可判定完成。

系统中共有两个班组，每个班组配备若干台不同类型设备。设备在不同工序之间可重复使用，但同一时刻每台设备只能服务一道工序。同一设备在同一车间内切换工序时不计运输时间，而跨车间作业时需要计入设备转运时间。转运时间由车间距离和设备移动速度确定，设备移动速度为 $2\,\mathrm{m/s}$，持续作业时间和运输时间均按秒向上取整。

题目要求解决以下四个递进问题：
\begin{enumerate}[label=\textbf{问题\arabic*.}, leftmargin=2.8em]
  \item 仅由班组 1 独立承担 A 车间全部整修任务，求最短完工时长并给出设备调度表。
  \item 仅使用班组 1 设备完成 A--E 五个车间全部任务，求最短完工时长并给出设备调度表。
  \item 使用班组 1 和班组 2 设备共同完成五个车间全部任务，求最短完工时长并给出带班组信息的设备调度表。
  \item 在 50 万元设备购置预算下，综合确定设备购置方案与作业调度策略，使全部任务完成时间最短，并给出调度表和采购表。
\end{enumerate}

该问题本质上是一个带运输时间的多资源调度问题。其难点主要体现在：车间内部工序必须严格按链式顺序执行；同类设备存在多台可替代机器；双设备工序需要同步协调；同一设备跨车间作业时存在序列相关转运时间；最后一问还引入了预算约束下的设备扩容决策。

从调度理论角度看，该问题兼具柔性作业车间调度和资源受限项目调度特征；同设备跨车间运输时间可视为序列相关准备时间。混合整数规划是小规模作业车间调度中常用的精确建模方法\cite{KuBeck2016}，而带序列相关准备时间的调度生成与排序约束也已有较成熟研究\cite{Artigues2005}。本文在完整 MILP 调度模型基础上，引入 CP-SAT 进行关键下界和采购模型验证，主要参考了近年来约束规划在调度问题中的应用思路\cite{PyJobShop2025,ORToolsCPSAT} 以及 SciPy/HiGHS 对 MILP 的求解接口\cite{SciPyMilp}。
    % 一、问题重述
% 02_analysis.tex —— 问题分析
\section{问题分析}

\subsection{工序展开与设备子任务拆分}

题目中每个车间的整修任务均构成一条工序链。例如 A 车间必须按 A1、A2、A3 顺序执行；D 车间必须按 D1 至 D6 顺序执行。C 车间较为特殊，除 C1、C2 外，C3、C4、C5 需要重复执行三轮，因此展开后 C 车间工序链为
\begin{align*}
\text{C1}&\rightarrow \text{C2}
\rightarrow(\text{C3\_R1}\rightarrow \text{C4\_R1}\rightarrow \text{C5\_R1})\\
&\rightarrow(\text{C3\_R2}\rightarrow \text{C4\_R2}\rightarrow \text{C5\_R2})\\
&\rightarrow(\text{C3\_R3}\rightarrow \text{C4\_R3}\rightarrow \text{C5\_R3}) .
\end{align*}
展开后，五个车间共形成 27 个工序实例。由于部分工序需要两类设备，各工序进一步拆分为 41 个设备子任务。设备子任务是模型中进行设备分配和互斥约束的基本单位。

\subsection{资源瓶颈与关键路径}

班组配置中，自动化输送臂、工业清洗机、精密灌装机数量较多，而高速抛光机和自动传感多功能机每个班组仅各 1 台。对于单班组场景，后两类设备天然成为瓶颈资源。问题二的最优性证明也正是围绕这两台单机瓶颈展开。

当两个班组共同参与后，设备数量翻倍，单台设备瓶颈有所缓解。此时 C 车间的固定工序链成为不可压缩的关键路径。由于车间内部工序必须严格顺序执行，新增设备只能改善不同车间或不同工序之间的并行能力，无法缩短 C 车间链上每一道工序的固有作业时长。因此，问题三和问题四的最短完工时间均由 C 车间关键链决定。

\subsection{四问递进关系}

四个问题的约束范围逐步扩大。问题一是单车间链式调度，可直接由工序链下界判断最优；问题二增加五车间并行和跨车间运输，核心矛盾变为单班组瓶颈设备排序；问题三增加第二班组，资源扩容后关键瓶颈从单机排序转为 C 车间不可压缩链；问题四进一步增加采购决策，但由于问题三已经达到 C 链下界，预算不再影响最优工期。

\begin{table}[!htbp]
  \centering
  \caption{四个问题最短时长汇总}
  \label{tab:summary}
  \renewcommand{\arraystretch}{1.25}
  \begin{tabularx}{\textwidth}{c X r c X}
    \toprule
    问题 & 资源条件 & 最短时长(s) & 时:分:秒 & 最优性依据 \\
    \midrule
    Q1 & 班组 1，A 车间 & 41600 & 11:33:20 & A 车间工序链下界 \\
    Q2 & 班组 1，五车间 & 163764 & 45:29:24 & 双瓶颈资源松弛下界 \\
    Q3 & 班组 1 + 2，五车间 & 123844 & 34:24:04 & C 车间关键链下界 \\
    Q4 & 双班组 + 50 万元预算 & 123844 & 34:24:04 & C 链下界 + 显式采购模型 \\
    \bottomrule
  \end{tabularx}
\end{table}

\begin{figure}[!htbp]
  \centering
  \includegraphics[width=0.86\textwidth]{makespan_summary.png}
  \caption{四问最短工期对比}
  \label{fig:makespan-summary}
\end{figure}
       % 二、问题分析
% 03_assumptions.tex —— 模型假设
\section{模型假设}

为使模型可计算且与题意一致，本文作如下假设：
\begin{enumerate}
  \item 每台设备初始位置为其所属班组驻地，首次进入某车间作业前必须满足从班组驻地到该车间的转运时间。
  \item 同一台设备在同一车间内连续执行不同工序时，设备转运时间忽略。
  \item 同一台设备跨车间作业时，转运时间为车间间距离除以 $2\,\mathrm{m/s}$ 后向上取整。
  \item 对需要两类设备的工序，两类设备同时投入并分别完成该工序对应工程量；工序完成时间取两类设备完成时间的较晚者。
  \item 所有设备子任务的工作时间均按秒计，并由工程量和设备效率计算后向上取整。
  \item 问题四中，新增购置设备与同类型原有设备具有相同效率、移动速度和使用规则，且购买后立即可用于调度。
  \item 优化目标首先为最小化总完工时间；在总完工时间相同情况下，问题四进一步选择采购费用最小的方案。
\end{enumerate}

其中假设 4 对双设备协同工序的处理符合题目``不考虑两类设备先后顺序、不考虑两类设备之间等待时间''的表述。并且在问题三、四中，调度方案已经达到与该假设无关的 C 车间关键链下界，因此该处理不会影响最终最短工期结论。
    % 三、模型假设
% 04_notation.tex —— 符号说明
\section{符号说明}
% 表格用 tabularx{\textwidth} 撑满版心宽，「含义」列用 X 自适应拉伸。
\begin{table}[!htbp]
  \centering
  \caption{符号说明}%
  \label{tab:notation}
  \renewcommand{\arraystretch}{1.35}
  \begin{tabularx}{\textwidth}{c X c}
    \toprule
    符号 & 含义 & 单位 \\
    \midrule
    $i$        & 工序编号                              & —   \\
    $m$        & 设备类型编号                          & —   \\
    $p_{im}$   & 工序 $i$ 在设备类型 $m$ 上的加工工时  & 秒  \\
    % TODO: 补全全部符号
    \bottomrule
  \end{tabularx}
\end{table}
       % 四、符号说明
% 05_model.tex —— 模型建立与求解
\section{模型建立与求解}

\subsection{数据预处理与工时计算}

对于设备子任务 $o$，其作业时长由工程量和设备效率计算：
\begin{equation}
  p_o=\left\lceil \frac{q_{j(o)}}{v_o}\times 3600 \right\rceil ,
  \label{eq:duration}
\end{equation}
其中 $q_{j(o)}$ 为工序工程量，$v_o$ 为对应设备效率，单位为 $\mathrm{m^3/h}$。若工序 $j$ 需要多类设备共同完成，则该工序链上不可压缩占用时间为
\begin{equation}
  d_j=\max_{o:j(o)=j} p_o .
  \label{eq:job-duration}
\end{equation}
例如，A1 需要精密灌装机和自动化输送臂。二者对应持续时间分别为 $5400\,\mathrm{s}$ 和 $4320\,\mathrm{s}$，因此 A1 的工序链占用时长为 $5400\,\mathrm{s}$。

若设备从位置 $u$ 移动到位置 $v$，转运时间为
\begin{equation}
  r_{u,v}=\left\lceil \frac{D_{u,v}}{2}\right\rceil ,
  \label{eq:travel}
\end{equation}
其中 $D_{u,v}$ 为题目附件给出的距离，单位为米。若 $u=v$，则 $r_{u,v}=0$。

附件给出的两个班组设备配置完全相同。各类型设备数量、速度和单价如表~\ref{tab:equipment} 所示。可以看出，高速抛光机和自动传感多功能机每个班组均只有 1 台，是后续调度中的天然瓶颈资源。

\begin{table}[!htbp]
  \centering
  \caption{班组设备配置表}
  \label{tab:equipment}
  \renewcommand{\arraystretch}{1.2}
  \begin{tabularx}{\textwidth}{c X r r r}
    \toprule
    班组 & 设备类型 & 设备数量 & 速度(m/s) & 单价(元) \\
    \midrule
    1 & 自动化输送臂 & 4 & 2 & 50000 \\
    1 & 工业清洗机 & 5 & 2 & 40000 \\
    1 & 精密灌装机 & 5 & 2 & 35000 \\
    1 & 自动传感多功能机 & 1 & 2 & 80000 \\
    1 & 高速抛光机 & 1 & 2 & 75000 \\
    2 & 自动化输送臂 & 4 & 2 & 50000 \\
    2 & 工业清洗机 & 5 & 2 & 40000 \\
    2 & 精密灌装机 & 5 & 2 & 35000 \\
    2 & 自动传感多功能机 & 1 & 2 & 80000 \\
    2 & 高速抛光机 & 1 & 2 & 75000 \\
    \bottomrule
  \end{tabularx}
\end{table}

展开后的主要工序时长如表~\ref{tab:process-duration} 所示。对于双设备工序，表中工序链占用时间为对应设备子任务时长的最大值。

\begin{table}[!htbp]
  \centering
  \caption{展开后主要工序时长摘要}
  \label{tab:process-duration}
  \small
  \renewcommand{\arraystretch}{1.18}
  \begin{tabularx}{\textwidth}{c c r X r}
    \toprule
    车间 & 工序 & 工程量($\mathrm{m^3}$) & 所需设备 & 工序链占用时间(s) \\
    \midrule
    A & A1 & 300 & 精密灌装机、自动化输送臂 & 5400 \\
    A & A2 & 500 & 高速抛光机、工业清洗机 & 18000 \\
    A & A3 & 500 & 自动传感多功能机 & 18000 \\
    B & B1 & 120 & 工业清洗机 & 4320 \\
    B & B2 & 1500 & 精密灌装机、自动化输送臂 & 27000 \\
    B & B3 & 360 & 精密灌装机 & 3703 \\
    B & B4 & 360 & 高速抛光机、自动传感多功能机 & 12960 \\
    C & C1 & 720 & 工业清洗机、自动化输送臂 & 10368 \\
    C & C2 & 720 & 精密灌装机 & 7406 \\
    C & C3\_Ri & 360 & 精密灌装机、自动化输送臂 & 6480 \\
    C & C4\_Ri & 400 & 高速抛光机、工业清洗机 & 14400 \\
    C & C5\_Ri & 400 & 自动传感多功能机 & 14400 \\
    D & D1 & 600 & 工业清洗机 & 8640 \\
    D & D2 & 800 & 精密灌装机、自动化输送臂 & 14400 \\
    D & D3 & 450 & 精密灌装机 & 4629 \\
    D & D4 & 1500 & 高速抛光机、自动传感多功能机 & 45000 \\
    D & D5 & 1500 & 自动传感多功能机 & 18000 \\
    D & D6 & 700 & 高速抛光机 & 25200 \\
    E & E1 & 1000 & 工业清洗机 & 14400 \\
    E & E2 & 600 & 精密灌装机 & 6172 \\
    E & E3 & 600 & 自动传感多功能机、工业清洗机 & 21600 \\
    \bottomrule
  \end{tabularx}
\end{table}

\subsection{固定设备集合下的混合整数规划模型}

问题一至问题三均可视为固定设备集合下的调度问题，区别仅在于允许使用的车间和设备集合不同。问题一只考虑 A 车间和班组 1；问题二考虑五个车间和班组 1；问题三考虑五个车间和两个班组。

目标函数为最小化全部任务完工时间：
\begin{equation}
  \min T .
  \label{eq:objective}
\end{equation}

每个设备子任务必须且只能分配给一台对应类型设备：
\begin{equation}
  \sum_{m\in M:\,\mathit{type}(m)=\tau(o)} x_{o,m}=1,\quad \forall o\in O .
  \label{eq:assignment}
\end{equation}

若工序 $a$ 和 $b$ 属于同一车间，且 $b$ 是 $a$ 的后续相邻工序，则
\begin{equation}
  S_b\ge S_a+d_a .
  \label{eq:precedence}
\end{equation}
该约束保证各车间内部工序严格按照题目给定顺序执行。对于任意工序 $j$，总完工时间不得早于该工序完成时间：
\begin{equation}
  T\ge S_j+d_j,\quad \forall j\in J .
  \label{eq:makespan}
\end{equation}

若子任务 $o$ 分配给设备 $m$，则所属工序开始时间不得早于设备从班组驻地到达该车间的时间：
\begin{equation}
  S_{j(o)}\ge r_{\mathit{origin}(m),w(j(o))}x_{o,m}.
  \label{eq:initial-arrival}
\end{equation}

若两个子任务 $a,b$ 由同一台设备 $m$ 执行，则二者不能重叠。引入二元排序变量 $y_{a,b,m}$，当 $x_{a,m}=x_{b,m}=1$ 且 $y_{a,b,m}=1$ 时，表示子任务 $a$ 在同一设备 $m$ 上排在 $b$ 之前；当 $y_{a,b,m}=0$ 时，表示 $b$ 排在 $a$ 之前。对应 big-M 线性化为
\begin{align}
S_{j(b)}-S_{j(a)}
&\ge p_a+r_{w(j(a)),w(j(b))}
-M(1-y_{a,b,m})-M(2-x_{a,m}-x_{b,m}), \label{eq:machine-order-1}\\
S_{j(a)}-S_{j(b)}
&\ge p_b+r_{w(j(b)),w(j(a))}
-My_{a,b,m}-M(2-x_{a,m}-x_{b,m}). \label{eq:machine-order-2}
\end{align}
当 $a,b$ 未同时分配给设备 $m$ 时，项 $M(2-x_{a,m}-x_{b,m})$ 使约束自动松弛；当二者同时分配给设备 $m$ 时，两条约束中恰有一条根据 $y_{a,b,m}$ 生效，从而确定二者的先后顺序。本文取 $M=10^6\,\mathrm{s}$，大于所有工序作业时间、转运时间与总工期的可能上界。

\subsection{求解流程与计算设置}

本文采用``精确求解 + 独立下界验证''的两层流程：
\begin{enumerate}
  \item 对问题一至问题三，建立固定设备集合下的 MILP 模型，以最小化完工时间 $T$ 为目标求解，得到可行调度表和最优目标值。
  \item 在得到最优完工时间后，增加约束 $T\le T^\ast$，再以工序开始时间之和为次级目标进行紧凑化，使非关键任务尽量提前，便于输出更紧凑的调度表。
  \item 对问题二，额外构造仅保留高速抛光机和自动传感多功能机两类单台瓶颈设备冲突的 CP-SAT 松弛模型，得到独立下界。
  \item 对问题四，构造含采购变量的显式 CP-SAT 模型，先最小化完工时间，再在最优完工时间下最小化采购费用。
\end{enumerate}

固定设备 MILP 使用 HiGHS 混合整数规划求解器。模型中的连续变量为工序开始时间和完工时间，二元变量为设备分配变量和同设备排序变量。CP-SAT 验证模型用于两个目的：一是对问题二构造瓶颈资源松弛下界；二是对问题四构造采购--调度联合模型。固定设备 MILP 的求解统计如表~\ref{tab:milp-stats} 所示。

\begin{table}[!htbp]
  \centering
  \caption{固定设备 MILP 求解统计}
  \label{tab:milp-stats}
  \renewcommand{\arraystretch}{1.2}
  \begin{tabularx}{\textwidth}{X r c r r r}
    \toprule
    子问题 & 最短时长(s) & 求解状态 & 节点数 & LP 迭代 & 求解时间(s) \\
    \midrule
    Q1 & 41600 & Optimal & 0 & 15 & 0.01 \\
    Q2 & 163764 & Optimal & 206 & 2600 & 0.46 \\
    Q3 & 123844 & Optimal & 800 & 3263 & 1.19 \\
    Q2 二阶段紧凑化 & 163764 & Optimal & 445 & 3382 & 0.45 \\
    Q3 二阶段紧凑化 & 123844 & Optimal & 1211 & 16386 & 2.15 \\
    \bottomrule
  \end{tabularx}
\end{table}

\subsection{问题一：班组 1 完成 A 车间}

问题一只涉及 A 车间，工序顺序为 $\text{A1}\rightarrow\text{A2}\rightarrow\text{A3}$。班组 1 到 A 车间距离为 $400\,\mathrm{m}$，设备速度为 $2\,\mathrm{m/s}$，因此最早到达时间为 $200\,\mathrm{s}$。A 车间三道工序的链式占用时间分别为 $5400\,\mathrm{s}$、$18000\,\mathrm{s}$ 和 $18000\,\mathrm{s}$。因此任意可行方案的完工时间均满足
\begin{equation}
  T\ge 200+5400+18000+18000=41600\,\mathrm{s}.
  \label{eq:q1-lb}
\end{equation}
优化模型给出的排程恰好达到该下界，故问题一最短完工时间为
\[
  \boxed{41600\,\mathrm{s}=11:33:20}.
\]

\begin{figure}[!htbp]
  \centering
  \includegraphics[width=\textwidth]{gantt_q1.png}
  \caption{问题一调度甘特图}
  \label{fig:q1-gantt}
\end{figure}

\begin{table}[!htbp]
  \centering
  \caption{问题一设备调度结果}
  \label{tab:q1-schedule}
  \renewcommand{\arraystretch}{1.18}
  \begin{tabularx}{\textwidth}{r X c c r c}
    \toprule
    序号 & 设备编号 & 起始时间 & 结束时间 & 持续工作时间(s) & 工序编号 \\
    \midrule
    1 & 精密灌装机1-1 & 00:03:20 & 01:33:20 & 5400 & A1 \\
    2 & 自动化输送臂1-4 & 00:03:20 & 01:15:20 & 4320 & A1 \\
    3 & 工业清洗机1-5 & 01:33:20 & 03:33:20 & 7200 & A2 \\
    4 & 高速抛光机1-1 & 01:33:20 & 06:33:20 & 18000 & A2 \\
    5 & 自动传感多功能机1-1 & 06:33:20 & 11:33:20 & 18000 & A3 \\
    \bottomrule
  \end{tabularx}
\end{table}

\subsection{问题二：班组 1 完成五车间}

问题二使用班组 1 的设备完成五个车间全部任务。相比问题一，问题二新增了多车间并行、设备跨车间转运和单班组资源竞争。班组 1 中高速抛光机和自动传感多功能机均只有 1 台。由于 A、B、C、D 等车间的关键工序均需要这两类设备，它们构成问题二的主要瓶颈。

为证明最优性，本文构造如下松弛模型：保留所有车间内部工序链约束、所有工序最短占用时间 $d_j$、班组 1 初始到达各车间的时间约束；仅对高速抛光机和自动传感多功能机保留设备互斥和跨车间转运时间；放松自动化输送臂、工业清洗机、精密灌装机的设备互斥冲突。该松弛模型删除了部分资源冲突，故其最优值是原问题的合法下界。CP-SAT 求解该松弛模型得到
\[
  LB_{Q2}=163764\,\mathrm{s}.
\]
完整调度模型同时给出一组完工时间为 $163764\,\mathrm{s}$ 的可行排程，因此
\[
  163764\le OPT(Q2)\le 163764.
\]
故在本文采用的双设备同步协同模型下，问题二全局最优完工时间为
\[
  \boxed{163764\,\mathrm{s}=45:29:24}.
\]

\begin{figure}[!htbp]
  \centering
  \includegraphics[width=\textwidth]{gantt_q2.png}
  \caption{问题二调度甘特图}
  \label{fig:q2-gantt}
\end{figure}

双瓶颈资源松弛模型中的关键设备序列如表~\ref{tab:q2-polishing} 和表~\ref{tab:q2-sensing} 所示。它们说明在单班组情况下，完工时间主要由高速抛光机和自动传感多功能机的占用顺序决定。完整的问题二设备调度表见附录表~\ref{tab:appendix-q2}。

\begin{table}[!htbp]
  \centering
  \caption{问题二高速抛光机瓶颈序列}
  \label{tab:q2-polishing}
  \renewcommand{\arraystretch}{1.15}
  \begin{tabularx}{\textwidth}{c c c c r}
    \toprule
    工序 & 车间 & 开始 & 结束 & 本设备时长(s) \\
    \midrule
    A2 & A & 01:33:20 & 06:33:20 & 18000 \\
    C4\_R1 & C & 06:48:04 & 10:08:04 & 12000 \\
    D4 & D & 10:12:24 & 22:42:24 & 45000 \\
    C4\_R2 & C & 22:46:44 & 26:06:44 & 12000 \\
    B4 & B & 31:55:54 & 34:55:54 & 10800 \\
    C4\_R3 & C & 35:05:04 & 38:25:04 & 12000 \\
    D6 & D & 38:29:24 & 45:29:24 & 25200 \\
    \bottomrule
  \end{tabularx}
\end{table}

\begin{table}[!htbp]
  \centering
  \caption{问题二自动传感多功能机瓶颈序列}
  \label{tab:q2-sensing}
  \renewcommand{\arraystretch}{1.15}
  \begin{tabularx}{\textwidth}{c c c c r}
    \toprule
    工序 & 车间 & 开始 & 结束 & 本设备时长(s) \\
    \midrule
    E3 & E & 05:46:12 & 07:46:12 & 7200 \\
    D4 & D & 10:12:24 & 15:12:24 & 18000 \\
    C5\_R1 & C & 15:16:44 & 19:16:44 & 14400 \\
    D5 & D & 22:42:24 & 27:42:24 & 18000 \\
    C5\_R2 & C & 27:46:44 & 31:46:44 & 14400 \\
    B4 & B & 31:55:54 & 35:31:54 & 12960 \\
    A3 & A & 35:40:24 & 40:40:24 & 18000 \\
    C5\_R3 & C & 40:49:09 & 44:49:09 & 14400 \\
    \bottomrule
  \end{tabularx}
\end{table}

\subsection{问题三：两班组共同完成五车间}

问题三允许使用班组 1 和班组 2 的全部设备。两班组设备配置相同，因此各类设备数量相比问题二翻倍。设备数量增加后，单台设备瓶颈显著缓解，但 C 车间工序链成为不可压缩的全局关键路径。

C 车间展开后的工序链为 $\text{C1}\rightarrow \text{C2}\rightarrow(\text{C3}\rightarrow\text{C4}\rightarrow\text{C5})\times 3$。其不含初始转运的最短链长为
\begin{align}
  L_C
  &=10368+7406+3(6480+14400+14400) \nonumber\\
  &=123614\,\mathrm{s}.
  \label{eq:c-chain}
\end{align}
两班组到 C 车间的初始转运时间分别为 $230\,\mathrm{s}$ 和 $310\,\mathrm{s}$，因此任意调度方案的完工时间满足
\begin{equation}
  T\ge 123614+230=123844\,\mathrm{s}.
  \label{eq:q3-lb}
\end{equation}
完整调度模型给出的两班组排程达到该下界，因此问题三全局最优完工时间为
\[
  \boxed{123844\,\mathrm{s}=34:24:04}.
\]

\begin{figure}[!htbp]
  \centering
  \includegraphics[width=\textwidth]{gantt_q3.png}
  \caption{问题三调度甘特图}
  \label{fig:q3-gantt}
\end{figure}

从图~\ref{fig:q3-gantt} 可以看出，在两班组共同参与后，多个车间工序能够并行推进，但 C 车间链上工序仍必须依次完成。最终完工时刻对应 C5\_R3 的完成，说明 C 车间关键链确实决定了全局工期。问题三完整设备调度表见附录表~\ref{tab:appendix-q3}。

\subsection{问题四：预算约束下的设备购置与调度}

问题四在问题三基础上增加 50 万元设备购置预算。对每一台候选增购设备 $m$，定义二元变量
\[
z_m=
\begin{cases}
1,&\text{购买设备 }m,\\
0,&\text{不购买设备 }m.
\end{cases}
\]
预算约束为
\begin{equation}
  \sum_m \mathit{price}_m z_m\le 500000 .
  \label{eq:budget}
\end{equation}
若某设备子任务分配给增购设备 $m$，则必须购买该设备：
\begin{equation}
  x_{o,m}\le z_m .
  \label{eq:purchase-activation}
\end{equation}
在此基础上，调度部分沿用固定设备模型中的工序链、设备互斥、初始到达和跨车间运输约束。

问题四采用字典序优化。第一阶段求解 $\min T$；在得到最优完工时间 $T^\ast$ 后，第二阶段固定 $T=T^\ast$，再最小化采购费用 $\sum_m \mathit{price}_m z_m$。显式采购模型求得第一阶段最优完工时间仍为
\[
  T^\ast=123844\,\mathrm{s}.
\]
第二阶段在固定 $T=123844\,\mathrm{s}$ 的条件下，得到最小采购费用为 0 元。因此问题四最优采购方案为所有设备均购买 0 台。

该结论也可由 C 车间关键链解释。问题三已经达到 C 车间关键链下界
\[
  123614+230=123844\,\mathrm{s}.
\]
新增设备只能增加不同工序或不同车间之间的并行能力，不能缩短 C 车间链上任一单工序的作业时长，也不能改变 C 车间内部严格顺序。因此任何采购方案均无法使完工时间低于 $123844\,\mathrm{s}$。

\begin{figure}[!htbp]
  \centering
  \includegraphics[width=\textwidth]{gantt_q4.png}
  \caption{问题四调度甘特图}
  \label{fig:q4-gantt}
\end{figure}

问题四设备购买情况如表~\ref{tab:q4-purchase} 所示。由于最优采购为 0 台，设备集合与问题三相同，最优调度策略可直接采用问题三的双班组调度表；问题四完整调度表见附录表~\ref{tab:appendix-q4}。

\begin{table}[!htbp]
  \centering
  \caption{问题四设备购买情况}
  \label{tab:q4-purchase}
  \renewcommand{\arraystretch}{1.2}
  \begin{tabularx}{\textwidth}{X r r r r}
    \toprule
    设备名称 & 班组1购买台数 & 班组2购买台数 & 设备单价(元) & 小计(元) \\
    \midrule
    自动化输送臂 & 0 & 0 & 50000 & 0 \\
    工业清洗机 & 0 & 0 & 40000 & 0 \\
    精密灌装机 & 0 & 0 & 35000 & 0 \\
    自动传感多功能机 & 0 & 0 & 80000 & 0 \\
    高速抛光机 & 0 & 0 & 75000 & 0 \\
    \bottomrule
  \end{tabularx}
\end{table}

为分析追加预算对最优工期的影响，考虑预算上限 $B$ 从 0 增加到 500000 元的情形。由于问题三在无新增设备时已经达到 C 车间关键链下界，且采购设备不能改变 C 车间内部严格顺序，也不能缩短单个工序的固有作业时长，因此任意非负预算下的最优工期均不可能低于该下界。预算敏感性结果如表~\ref{tab:budget-sensitivity} 所示。

\begin{table}[!htbp]
  \centering
  \caption{预算敏感性分析}
  \label{tab:budget-sensitivity}
  \renewcommand{\arraystretch}{1.2}
  \begin{tabularx}{\textwidth}{r X r c c}
    \toprule
    预算上限(元) & 可行采购方案 & 最优完工时间(s) & HH:MM:SS & 是否达到 C 链下界 \\
    \midrule
    0 & 不采购 & 123844 & 34:24:04 & 是 \\
    100000 & 不采购 & 123844 & 34:24:04 & 是 \\
    200000 & 不采购 & 123844 & 34:24:04 & 是 \\
    300000 & 不采购 & 123844 & 34:24:04 & 是 \\
    400000 & 不采购 & 123844 & 34:24:04 & 是 \\
    500000 & 不采购 & 123844 & 34:24:04 & 是 \\
    \bottomrule
  \end{tabularx}
\end{table}
          % 五、模型建立与求解（Q1–Q4）
% 06_evaluation.tex —— 模型评价与推广
\section{模型评价与推广}
\subsection{模型优点}
% TODO
\subsection{模型缺点}
% TODO
\subsection{模型推广}
% TODO: 对更一般 FJSP/RCPSP+SDST 实例的适用性与规模边界。
     % 六、模型评价与推广

\clearpage
\printbibliography                  % 七、参考文献

\clearpage
\appendix
% 07_appendix.tex —— 附录
\section{程序代码}

本文的完整求解、验证和绘图脚本保存在工程目录中。本附录列出关键代码片段，完整脚本路径如下：
\begin{itemize}
  \item \path{baseline_v1/code/model_solver.py}：MILP 建模求解并导出四问调度表；
  \item \path{code/analysis/verify_optimality_chain.py}：Q2 下界与 Q4 采购模型验证；
  \item \path{code/plot/plot_gantt.py}：根据 CSV 结果表重新生成论文图件。
\end{itemize}

\lstset{basicstyle=\scriptsize\ttfamily, breakatwhitespace=false}

\subsection{MILP 调度模型核心代码}
\lstinputlisting[
  language=Python,
  firstline=165,
  lastline=284,
  caption={MILP 变量、目标与约束构造核心代码}
]{../../baseline_v1/code/model_solver.py}

\subsection{结果导出与四问求解入口}
\lstinputlisting[
  language=Python,
  firstline=326,
  lastline=386,
  caption={四问求解流程与结果导出代码}
]{../../baseline_v1/code/model_solver.py}

\subsection{CP-SAT 最优性验证核心代码}
\lstinputlisting[
  language=Python,
  firstline=69,
  lastline=134,
  caption={问题二双瓶颈资源松弛下界模型}
]{../../code/analysis/verify_optimality_chain.py}

\lstinputlisting[
  language=Python,
  firstline=269,
  lastline=318,
  caption={问题四显式采购模型的两阶段优化代码}
]{../../code/analysis/verify_optimality_chain.py}

\subsection{论文图件生成代码}
\lstinputlisting[
  language=Python,
  firstline=53,
  lastline=145,
  caption={甘特图与工期汇总图生成代码}
]{../../code/plot/plot_gantt.py}

\section{补充图表}

本节给出正文未完整展开的题目要求结果表。表格均由 \path{baseline_v1/results/} 中的 CSV 文件自动生成：问题二、三、四的完整设备调度表分别对应 \path{table2_question2.csv}、\path{table3_question3.csv}、\path{table4_question4.csv}，设备购置表对应 \path{table5_purchase_plan.csv}，四问汇总表对应 \path{summary.csv}。

% This file is generated from baseline_v1/results/*.csv. Do not edit numbers by hand.

\begingroup
\scriptsize
\renewcommand{\arraystretch}{1.08}
\begin{longtable}{@{}r p{3.6cm} c c r c@{}}
\caption{问题二完整设备调度表}\label{tab:appendix-q2}\\
\toprule
序号 & 设备编号 & 起始时间 & 结束时间 & 持续工作时间(s) & 工序编号 \\
\midrule
\endfirsthead
\multicolumn{6}{c}{续表：问题二完整设备调度表}\\
\toprule
序号 & 设备编号 & 起始时间 & 结束时间 & 持续工作时间(s) & 工序编号 \\
\midrule
\endhead
\bottomrule
\endfoot
1 & 精密灌装机1-4 & 00:03:20 & 01:33:20 & 5400 & A1 \\
2 & 自动化输送臂1-2 & 00:03:20 & 01:15:20 & 4320 & A1 \\
3 & 工业清洗机1-4 & 00:03:20 & 04:03:20 & 14400 & E1 \\
4 & 工业清洗机1-2 & 00:03:50 & 02:56:38 & 10368 & C1 \\
5 & 自动化输送臂1-3 & 00:03:50 & 02:56:38 & 10368 & C1 \\
6 & 工业清洗机1-5 & 00:05:10 & 01:17:10 & 4320 & B1 \\
7 & 工业清洗机1-3 & 00:05:55 & 02:29:55 & 8640 & D1 \\
8 & 精密灌装机1-2 & 01:17:10 & 08:47:10 & 27000 & B2 \\
9 & 自动化输送臂1-1 & 01:17:10 & 06:17:10 & 18000 & B2 \\
10 & 工业清洗机1-1 & 01:33:20 & 03:33:20 & 7200 & A2 \\
11 & 高速抛光机1-1 & 01:33:20 & 06:33:20 & 18000 & A2 \\
12 & 精密灌装机1-4 & 02:29:55 & 06:29:55 & 14400 & D2 \\
13 & 自动化输送臂1-4 & 02:29:55 & 05:09:55 & 9600 & D2 \\
14 & 精密灌装机1-5 & 02:56:38 & 05:00:04 & 7406 & C2 \\
15 & 精密灌装机1-1 & 04:03:20 & 05:46:12 & 6172 & E2 \\
16 & 精密灌装机1-3 & 05:00:04 & 06:48:04 & 6480 & C3\_R1 \\
17 & 自动化输送臂1-2 & 05:00:04 & 06:26:28 & 5184 & C3\_R1 \\
18 & 工业清洗机1-2 & 05:46:12 & 11:46:12 & 21600 & E3 \\
19 & 自动传感多功能机1-1 & 05:46:12 & 07:46:12 & 7200 & E3 \\
20 & 精密灌装机1-5 & 06:29:55 & 07:47:04 & 4629 & D3 \\
21 & 工业清洗机1-4 & 06:48:04 & 10:48:04 & 14400 & C4\_R1 \\
22 & 高速抛光机1-1 & 06:48:04 & 10:08:04 & 12000 & C4\_R1 \\
23 & 精密灌装机1-5 & 08:47:10 & 09:48:53 & 3703 & B3 \\
24 & 自动传感多功能机1-1 & 10:12:24 & 15:12:24 & 18000 & D4 \\
25 & 高速抛光机1-1 & 10:12:24 & 22:42:24 & 45000 & D4 \\
26 & 自动传感多功能机1-1 & 15:16:44 & 19:16:44 & 14400 & C5\_R1 \\
27 & 精密灌装机1-1 & 19:16:44 & 21:04:44 & 6480 & C3\_R2 \\
28 & 自动化输送臂1-1 & 19:16:44 & 20:43:08 & 5184 & C3\_R2 \\
29 & 自动传感多功能机1-1 & 22:42:24 & 27:42:24 & 18000 & D5 \\
30 & 工业清洗机1-3 & 22:46:44 & 26:46:44 & 14400 & C4\_R2 \\
31 & 高速抛光机1-1 & 22:46:44 & 26:06:44 & 12000 & C4\_R2 \\
32 & 自动传感多功能机1-1 & 27:46:44 & 31:46:44 & 14400 & C5\_R2 \\
33 & 精密灌装机1-1 & 31:46:44 & 33:34:44 & 6480 & C3\_R3 \\
34 & 自动化输送臂1-4 & 31:46:44 & 33:13:08 & 5184 & C3\_R3 \\
35 & 自动传感多功能机1-1 & 31:55:54 & 35:31:54 & 12960 & B4 \\
36 & 高速抛光机1-1 & 31:55:54 & 34:55:54 & 10800 & B4 \\
37 & 工业清洗机1-5 & 35:05:04 & 39:05:04 & 14400 & C4\_R3 \\
38 & 高速抛光机1-1 & 35:05:04 & 38:25:04 & 12000 & C4\_R3 \\
39 & 自动传感多功能机1-1 & 35:40:24 & 40:40:24 & 18000 & A3 \\
40 & 高速抛光机1-1 & 38:29:24 & 45:29:24 & 25200 & D6 \\
41 & 自动传感多功能机1-1 & 40:49:09 & 44:49:09 & 14400 & C5\_R3 \\
\end{longtable}
\endgroup

\begingroup
\scriptsize
\renewcommand{\arraystretch}{1.08}
\begin{longtable}{@{}r p{3.4cm} c c r c c@{}}
\caption{问题三完整设备调度表}\label{tab:appendix-q3}\\
\toprule
序号 & 设备编号 & 起始时间 & 结束时间 & 持续工作时间(s) & 工序编号 & 班组 \\
\midrule
\endfirsthead
\multicolumn{7}{c}{续表：问题三完整设备调度表}\\
\toprule
序号 & 设备编号 & 起始时间 & 结束时间 & 持续工作时间(s) & 工序编号 & 班组 \\
\midrule
\endhead
\bottomrule
\endfoot
1 & 精密灌装机1-1 & 00:03:20 & 01:33:20 & 5400 & A1 & 1 \\
2 & 自动化输送臂1-1 & 00:03:20 & 01:15:20 & 4320 & A1 & 1 \\
3 & 工业清洗机1-4 & 00:03:20 & 04:03:20 & 14400 & E1 & 1 \\
4 & 工业清洗机2-2 & 00:03:50 & 01:15:50 & 4320 & B1 & 2 \\
5 & 工业清洗机1-5 & 00:03:50 & 02:56:38 & 10368 & C1 & 1 \\
6 & 自动化输送臂1-2 & 00:03:50 & 02:56:38 & 10368 & C1 & 1 \\
7 & 工业清洗机2-4 & 00:05:40 & 02:29:40 & 8640 & D1 & 2 \\
8 & 精密灌装机2-2 & 01:15:50 & 08:45:50 & 27000 & B2 & 2 \\
9 & 自动化输送臂1-4 & 01:15:50 & 06:15:50 & 18000 & B2 & 1 \\
10 & 工业清洗机2-2 & 01:33:20 & 03:33:20 & 7200 & A2 & 2 \\
11 & 高速抛光机2-1 & 01:33:20 & 06:33:20 & 18000 & A2 & 2 \\
12 & 精密灌装机2-1 & 02:29:40 & 06:29:40 & 14400 & D2 & 2 \\
13 & 自动化输送臂2-2 & 02:29:40 & 05:09:40 & 9600 & D2 & 2 \\
14 & 精密灌装机2-5 & 02:56:38 & 05:00:04 & 7406 & C2 & 2 \\
15 & 精密灌装机1-3 & 04:03:20 & 05:46:12 & 6172 & E2 & 1 \\
16 & 精密灌装机1-5 & 05:00:04 & 06:48:04 & 6480 & C3\_R1 & 1 \\
17 & 自动化输送臂1-2 & 05:00:04 & 06:26:28 & 5184 & C3\_R1 & 1 \\
18 & 工业清洗机1-2 & 05:46:12 & 11:46:12 & 21600 & E3 & 1 \\
19 & 自动传感多功能机1-1 & 05:46:12 & 07:46:12 & 7200 & E3 & 1 \\
20 & 精密灌装机1-2 & 06:29:40 & 07:46:49 & 4629 & D3 & 1 \\
21 & 工业清洗机2-3 & 06:48:04 & 10:48:04 & 14400 & C4\_R1 & 2 \\
22 & 高速抛光机1-1 & 06:48:04 & 10:08:04 & 12000 & C4\_R1 & 1 \\
23 & 自动传感多功能机2-1 & 07:46:49 & 12:46:49 & 18000 & D4 & 2 \\
24 & 高速抛光机2-1 & 07:46:49 & 20:16:49 & 45000 & D4 & 2 \\
25 & 精密灌装机2-5 & 08:45:50 & 09:47:33 & 3703 & B3 & 2 \\
26 & 自动传感多功能机1-1 & 10:48:04 & 14:48:04 & 14400 & C5\_R1 & 1 \\
27 & 自动传感多功能机2-1 & 13:00:24 & 16:36:24 & 12960 & B4 & 2 \\
28 & 高速抛光机1-1 & 13:00:24 & 16:00:24 & 10800 & B4 & 1 \\
29 & 精密灌装机1-4 & 14:48:04 & 16:36:04 & 6480 & C3\_R2 & 1 \\
30 & 自动化输送臂2-3 & 14:48:04 & 16:14:28 & 5184 & C3\_R2 & 2 \\
31 & 自动传感多功能机1-1 & 14:56:49 & 19:56:49 & 18000 & A3 & 1 \\
32 & 工业清洗机1-3 & 16:36:04 & 20:36:04 & 14400 & C4\_R2 & 1 \\
33 & 高速抛光机1-1 & 16:36:04 & 19:56:04 & 12000 & C4\_R2 & 1 \\
34 & 自动传感多功能机2-1 & 20:16:49 & 25:16:49 & 18000 & D5 & 2 \\
35 & 自动传感多功能机1-1 & 20:36:04 & 24:36:04 & 14400 & C5\_R2 & 1 \\
36 & 精密灌装机1-5 & 24:36:04 & 26:24:04 & 6480 & C3\_R3 & 1 \\
37 & 自动化输送臂1-3 & 24:36:04 & 26:02:28 & 5184 & C3\_R3 & 1 \\
38 & 高速抛光机2-1 & 25:16:49 & 32:16:49 & 25200 & D6 & 2 \\
39 & 工业清洗机2-2 & 26:24:04 & 30:24:04 & 14400 & C4\_R3 & 2 \\
40 & 高速抛光机1-1 & 26:24:04 & 29:44:04 & 12000 & C4\_R3 & 1 \\
41 & 自动传感多功能机1-1 & 30:24:04 & 34:24:04 & 14400 & C5\_R3 & 1 \\
\end{longtable}
\endgroup

\begingroup
\scriptsize
\renewcommand{\arraystretch}{1.08}
\begin{longtable}{@{}r p{3.4cm} c c r c c@{}}
\caption{问题四完整设备调度表}\label{tab:appendix-q4}\\
\toprule
序号 & 设备编号 & 起始时间 & 结束时间 & 持续工作时间(s) & 工序编号 & 班组 \\
\midrule
\endfirsthead
\multicolumn{7}{c}{续表：问题四完整设备调度表}\\
\toprule
序号 & 设备编号 & 起始时间 & 结束时间 & 持续工作时间(s) & 工序编号 & 班组 \\
\midrule
\endhead
\bottomrule
\endfoot
1 & 精密灌装机1-1 & 00:03:20 & 01:33:20 & 5400 & A1 & 1 \\
2 & 自动化输送臂1-1 & 00:03:20 & 01:15:20 & 4320 & A1 & 1 \\
3 & 工业清洗机1-4 & 00:03:20 & 04:03:20 & 14400 & E1 & 1 \\
4 & 工业清洗机2-2 & 00:03:50 & 01:15:50 & 4320 & B1 & 2 \\
5 & 工业清洗机1-5 & 00:03:50 & 02:56:38 & 10368 & C1 & 1 \\
6 & 自动化输送臂1-2 & 00:03:50 & 02:56:38 & 10368 & C1 & 1 \\
7 & 工业清洗机2-4 & 00:05:40 & 02:29:40 & 8640 & D1 & 2 \\
8 & 精密灌装机2-2 & 01:15:50 & 08:45:50 & 27000 & B2 & 2 \\
9 & 自动化输送臂1-4 & 01:15:50 & 06:15:50 & 18000 & B2 & 1 \\
10 & 工业清洗机2-2 & 01:33:20 & 03:33:20 & 7200 & A2 & 2 \\
11 & 高速抛光机2-1 & 01:33:20 & 06:33:20 & 18000 & A2 & 2 \\
12 & 精密灌装机2-1 & 02:29:40 & 06:29:40 & 14400 & D2 & 2 \\
13 & 自动化输送臂2-2 & 02:29:40 & 05:09:40 & 9600 & D2 & 2 \\
14 & 精密灌装机2-5 & 02:56:38 & 05:00:04 & 7406 & C2 & 2 \\
15 & 精密灌装机1-3 & 04:03:20 & 05:46:12 & 6172 & E2 & 1 \\
16 & 精密灌装机1-5 & 05:00:04 & 06:48:04 & 6480 & C3\_R1 & 1 \\
17 & 自动化输送臂1-2 & 05:00:04 & 06:26:28 & 5184 & C3\_R1 & 1 \\
18 & 工业清洗机1-2 & 05:46:12 & 11:46:12 & 21600 & E3 & 1 \\
19 & 自动传感多功能机1-1 & 05:46:12 & 07:46:12 & 7200 & E3 & 1 \\
20 & 精密灌装机1-2 & 06:29:40 & 07:46:49 & 4629 & D3 & 1 \\
21 & 工业清洗机2-3 & 06:48:04 & 10:48:04 & 14400 & C4\_R1 & 2 \\
22 & 高速抛光机1-1 & 06:48:04 & 10:08:04 & 12000 & C4\_R1 & 1 \\
23 & 自动传感多功能机2-1 & 07:46:49 & 12:46:49 & 18000 & D4 & 2 \\
24 & 高速抛光机2-1 & 07:46:49 & 20:16:49 & 45000 & D4 & 2 \\
25 & 精密灌装机2-5 & 08:45:50 & 09:47:33 & 3703 & B3 & 2 \\
26 & 自动传感多功能机1-1 & 10:48:04 & 14:48:04 & 14400 & C5\_R1 & 1 \\
27 & 自动传感多功能机2-1 & 13:00:24 & 16:36:24 & 12960 & B4 & 2 \\
28 & 高速抛光机1-1 & 13:00:24 & 16:00:24 & 10800 & B4 & 1 \\
29 & 精密灌装机1-4 & 14:48:04 & 16:36:04 & 6480 & C3\_R2 & 1 \\
30 & 自动化输送臂2-3 & 14:48:04 & 16:14:28 & 5184 & C3\_R2 & 2 \\
31 & 自动传感多功能机1-1 & 14:56:49 & 19:56:49 & 18000 & A3 & 1 \\
32 & 工业清洗机1-3 & 16:36:04 & 20:36:04 & 14400 & C4\_R2 & 1 \\
33 & 高速抛光机1-1 & 16:36:04 & 19:56:04 & 12000 & C4\_R2 & 1 \\
34 & 自动传感多功能机2-1 & 20:16:49 & 25:16:49 & 18000 & D5 & 2 \\
35 & 自动传感多功能机1-1 & 20:36:04 & 24:36:04 & 14400 & C5\_R2 & 1 \\
36 & 精密灌装机1-5 & 24:36:04 & 26:24:04 & 6480 & C3\_R3 & 1 \\
37 & 自动化输送臂1-3 & 24:36:04 & 26:02:28 & 5184 & C3\_R3 & 1 \\
38 & 高速抛光机2-1 & 25:16:49 & 32:16:49 & 25200 & D6 & 2 \\
39 & 工业清洗机2-2 & 26:24:04 & 30:24:04 & 14400 & C4\_R3 & 2 \\
40 & 高速抛光机1-1 & 26:24:04 & 29:44:04 & 12000 & C4\_R3 & 1 \\
41 & 自动传感多功能机1-1 & 30:24:04 & 34:24:04 & 14400 & C5\_R3 & 1 \\
\end{longtable}
\endgroup

\begingroup
\scriptsize
\renewcommand{\arraystretch}{1.08}
\begin{longtable}{@{}p{4.0cm} r r r r@{}}
\caption{问题四设备购买表}\label{tab:appendix-purchase}\\
\toprule
设备名称 & 班组1购买台数 & 班组2购买台数 & 设备单价(元) & 小计(元) \\
\midrule
\endfirsthead
\multicolumn{5}{c}{续表：问题四设备购买表}\\
\toprule
设备名称 & 班组1购买台数 & 班组2购买台数 & 设备单价(元) & 小计(元) \\
\midrule
\endhead
\bottomrule
\endfoot
自动化输送臂 & 0 & 0 & 50000 & 0 \\
工业清洗机 & 0 & 0 & 40000 & 0 \\
精密灌装机 & 0 & 0 & 35000 & 0 \\
自动传感多功能机 & 0 & 0 & 80000 & 0 \\
高速抛光机 & 0 & 0 & 75000 & 0 \\
\end{longtable}
\endgroup

\begingroup
\scriptsize
\renewcommand{\arraystretch}{1.08}
\begin{longtable}{@{}c r c p{6.2cm}@{}}
\caption{四问最短时长汇总表}\label{tab:appendix-summary}\\
\toprule
问题 & 最短时长(s) & 最短时长(HH:MM:SS) & 求解状态 \\
\midrule
\endfirsthead
\multicolumn{4}{c}{续表：四问最短时长汇总表}\\
\toprule
问题 & 最短时长(s) & 最短时长(HH:MM:SS) & 求解状态 \\
\midrule
\endhead
\bottomrule
\endfoot
Q1 & 41600 & 11:33:20 & Optimization terminated successfully. (HiGHS Status 7: Optimal) \\
Q2 & 163764 & 45:29:24 & Optimization terminated successfully. (HiGHS Status 7: Optimal) \\
Q3 & 123844 & 34:24:04 & Optimization terminated successfully. (HiGHS Status 7: Optimal) \\
Q4 & 123844 & 34:24:04 & 关键路径下界已达到；最优采购为0 \\
\end{longtable}
\endgroup

       % 附录

\end{document}
